\section{Neural Networks and Backpropagation}
\label{sec:neural-networks}

Neural networks are one way to replace manually designed decision rules by
learned functions. For image analysis, the input may be an image, a patch, or a
feature vector, and the output may be a class score, a probability, a
segmentation label, or another task-specific prediction. The central idea is
simple: compose many parameterized transformations and fit their parameters
from examples.

Szeliski's treatment of deep learning, especially the part on deep neural
networks and backpropagation, is a useful companion for this topic. The focus
here is narrower: understand the computational structure of a feed-forward
network, why backpropagation is the efficient way to compute the gradients
needed for training, and how to tell whether a trained network will generalize.

\subsection{From Perceptrons to Multilayer Networks}

A perceptron is a linear classifier. Given an input vector $x$, weights $w$,
and bias $b$, it predicts a binary label by thresholding an affine score:
\begin{equation}
  f(x) =
  \begin{cases}
    1, & w^\top x + b > 0,\\
    -1, & \text{otherwise}.
  \end{cases}
\end{equation}
\noindent
The decision boundary is the hyperplane $w^\top x + b = 0$. This makes the
perceptron easy to interpret, but also limited: if the classes cannot be
separated by a linear boundary, the standard perceptron cannot solve the
problem. The XOR problem is the usual minimal example.

During perceptron training, an incorrectly predicted example changes the
weights in the direction of the input:
\begin{equation}
  w_i^{t+1} = w_i^t + \eta\left(y - f(x)\right)x_i ,
\end{equation}
\noindent
where $\eta$ is the learning rate. The rule is useful historically because it
already shows the pattern of learning from prediction error, but modern neural
networks need a smoother and more general gradient-based view.

Logistic regression is the smooth relative of the perceptron and a useful
reference point. For $C$ classes and a $D$-dimensional input it computes one
affine score per class, $z = Wx + b$ with $W \in \mathbb{R}^{C\times D}$ and
$b\in\mathbb{R}^C$, and turns the scores into probabilities with the softmax
function (introduced below). Its parameter count, $C(D+1)$, therefore depends on
both the number of classes and the input dimensionality (for two classes one
score suffices, giving $D+1$). Despite the nonlinear softmax, the predicted class
is the largest affine score, and the boundary between two classes, where
$z_i = z_j$, is a hyperplane: logistic regression is a \emph{linear} classifier,
exactly like linear regression used for classification.

A multilayer perceptron (MLP) overcomes the purely linear limitation by stacking
several affine maps and nonlinear activation functions. For layer $l$,
\begin{equation}
  h^l = W^l x^l + b^l,\qquad
  y^l = \phi(h^l),\qquad
  x^{l+1}=y^l .
\end{equation}
\noindent
Without the nonlinearity $\phi$, stacking layers would still collapse into one
linear transformation. The nonlinearities make it possible to build complex
decision boundaries from many simple units. The entries of $W^l$ are the
\emph{weights}, one per connection, and $b^l$ is the \emph{bias}. Network
diagrams often draw the bias as an extra neuron with constant output $1$ whose
outgoing weights are the bias values; both views compute the same affine map.

\begin{figure}[ht]
  \centering
  \begin{tikzpicture}[
      node distance=1.15cm,
      unit/.style={circle, draw, minimum size=0.52cm, inner sep=0pt},
      label/.style={align=center, font=\small},
      arrow/.style={-{Latex[length=2.2mm]}, thick}
    ]
    \node[unit, fill=green!8] (x1) {};
    \node[unit, fill=green!8, below=0.55cm of x1] (x2) {};
    \node[unit, fill=green!8, below=0.55cm of x2] (x3) {};

    \node[unit, fill=blue!8, right=1.5cm of x1, yshift=-0.28cm] (h11) {};
    \node[unit, fill=blue!8, below=0.55cm of h11] (h12) {};
    \node[unit, fill=blue!8, below=0.55cm of h12] (h13) {};

    \node[unit, fill=blue!8, right=1.5cm of h11] (h21) {};
    \node[unit, fill=blue!8, below=0.55cm of h21] (h22) {};
    \node[unit, fill=blue!8, below=0.55cm of h22] (h23) {};

    \node[unit, fill=orange!12, right=1.5cm of h22] (out) {};

    \foreach \a in {x1,x2,x3} {
      \foreach \b in {h11,h12,h13} {
        \draw[arrow, gray!70] (\a) -- (\b);
      }
    }
    \foreach \a in {h11,h12,h13} {
      \foreach \b in {h21,h22,h23} {
        \draw[arrow, gray!70] (\a) -- (\b);
      }
    }
    \foreach \a in {h21,h22,h23} {
      \draw[arrow, gray!70] (\a) -- (out);
    }

    \node[label, above=0.28cm of x1] {input};
    \node[label, above=0.28cm of h11] {hidden};
    \node[label, above=0.28cm of h21] {hidden};
    \node[label, above=0.28cm of out] {output};
  \end{tikzpicture}
  \caption{A feed-forward network composes affine transformations and
  nonlinearities. Training changes the weights on the connections.}
\end{figure}

\noindent
Beyond being nonlinear, a useful transfer function is typically monotonic and
smooth, so that it has a well-defined derivative everywhere; the classical
logistic sigmoid and hyperbolic tangent are additionally bounded, which keeps
activations in a fixed range. Common activation functions include the logistic
sigmoid, hyperbolic tangent, and rectified linear unit:
\begin{align}
  \sigma(z) &= \frac{1}{1 + e^{-z}},
  & \sigma'(z) &= \sigma(z)\left(1-\sigma(z)\right),\\
  \tanh'(z) &= 1-\tanh^2(z),
  & \operatorname{ReLU}(z) &= \max(0,z).
\end{align}
\noindent
For ReLU, the derivative is $1$ for positive inputs and $0$ for negative
inputs; at zero one chooses a convenient subgradient in implementations. This
simple derivative is one reason ReLU-style activations are computationally
attractive, and ReLU is the standard choice in convolutional networks.

\paragraph{Worked forward pass.}
Computing a small network by hand is the best check that the notation is
understood. Take two inputs, one hidden layer with two ReLU units, and two ReLU
output units, with the weights of Table~\ref{tab:mlp-example}. The input is
$x = (2,1)$ and the target is $y = (0,1)$.

\begin{table}[ht]
  \centering
  \begin{tabular}{@{}lccc@{}}
    \toprule
    Unit & weight from 1st input & weight from 2nd input & bias \\
    \midrule
    hidden $h_1$ & $1$ & $-1$ & $0$ \\
    hidden $h_2$ & $-1$ & $1$ & $0.5$ \\
    output $o_1$ (inputs $h_1,h_2$) & $1$ & $1$ & $-1$ \\
    output $o_2$ (inputs $h_1,h_2$) & $2$ & $-1$ & $0$ \\
    \bottomrule
  \end{tabular}
  \caption{Weights of a 2--2--2 ReLU network. Each row lists the weights of the
  connections entering one unit; the bias is the weight from the constant-$1$
  neuron.}
  \label{tab:mlp-example}
\end{table}

\noindent
Layer by layer, first form the weighted sum and then apply ReLU:
\begin{align*}
  h_1 &= \max(0,\ 1\cdot 2 - 1\cdot 1 + 0) = 1, &
  h_2 &= \max(0,\ -1\cdot 2 + 1\cdot 1 + 0.5) = \max(0,-0.5) = 0,\\
  o_1 &= \max(0,\ 1\cdot 1 + 1\cdot 0 - 1) = 0, &
  o_2 &= \max(0,\ 2\cdot 1 - 1\cdot 0 + 0) = 2 .
\end{align*}
\noindent
With the squared-error criterion $e = \tfrac12\lVert o - y\rVert_2^2$ the error
is $e = \tfrac12\bigl((0-0)^2 + (2-1)^2\bigr) = \tfrac12$. Two details are easy
to get wrong under time pressure: ReLU clips negative sums to zero \emph{before}
they enter the next layer, and the factor $\tfrac12$ belongs to the whole sum.

\subsection{Losses, Optimization, and the Need for Gradients}

Training starts with examples $(x_i,y_i)$ and a network $f_\theta$ with
parameters $\theta$, meaning all weights and biases. A loss function measures
how bad a prediction is. Over a dataset, the empirical objective is
\begin{equation}
  L(\theta) =
  \sum_{i=1}^{N} \ell\!\left(f_\theta(x_i), y_i\right).
\end{equation}
\noindent
For regression, $\ell$ is typically the squared error
$\tfrac12\lVert f_\theta(x) - y\rVert^2$. For classification, the network
outputs one score (logit) $z_c$ per class, the softmax turns the scores into
probabilities, and the loss is the cross-entropy, the negative log-probability
of the correct class $y$:
\begin{equation}
  p_c = \frac{e^{z_c}}{\sum_{j=1}^{C} e^{z_j}},
  \qquad
  \ell_{\mathrm{CE}} = -\log p_y .
\end{equation}
\noindent
The loss is what the optimizer sees, so it must match the task; accuracy itself
is not differentiable and is only used to evaluate the result.

Even when the loss is convex, as for logistic regression, there is in general no
closed-form solution for the best parameters, so the minimum has to be found
iteratively. Gradient descent is therefore not a tool reserved for non-convex
problems; convexity only guarantees that it will not get stuck in a worse local
minimum. For deep networks the loss is non-convex, and (stochastic) gradient
descent gives no guarantee of reaching the optimal solution.

Gradient descent steps against the gradient with learning rate $\eta$. The
three classical variants differ only in how many examples enter one update:
\begin{align}
  \text{online:}\quad
    &\Delta w = -\eta\, \frac{\partial e_n}{\partial w}
      &&\text{after every single example } n,\\
  \text{mini-batch:}\quad
    &\Delta w = -\eta \sum_{n\in\mathcal{B}_t} \frac{\partial e_n}{\partial w}
      &&\text{after a small batch } \mathcal{B}_t,\\
  \text{offline (batch):}\quad
    &\Delta w = -\eta \sum_{n=1}^{N} \frac{\partial e_n}{\partial w}
      &&\text{after a pass over all } N \text{ examples}.
\end{align}
\noindent
Offline learning follows the exact gradient of the full objective but updates
rarely; online learning updates often but with a noisy direction. Stochastic
gradient descent in modern practice is the mini-batch variant. Two common
additions modify the update. \emph{Momentum} (inertia) keeps part of the previous
step, so that consistent directions accelerate and oscillations cancel,
\begin{equation}
  \Delta w(t+1) = -\eta\, \frac{\partial e}{\partial w} + \alpha\, \Delta w(t),
  \qquad 0 \le \alpha < 1 ;
\end{equation}
\noindent
\emph{weight decay} adds a penalty $\tfrac{\lambda}{2}\lVert w\rVert^2$ to the
loss, which contributes $-\eta\lambda w$ to every step and pulls unused weights
toward zero.

Whatever the variant, every step needs the partial derivatives of the loss with
respect to all parameters. The practical difficulty is not writing the update
down; it is computing these derivatives for millions of parameters quickly and
without duplicating work. Backpropagation solves exactly that problem.

\begin{table}[ht]
  \centering
  \begin{tabular}{@{}p{3.0cm}p{5.1cm}p{5.1cm}@{}}
    \toprule
    Component & Role in the network & Training question \\
    \midrule
    Weights and biases & Store the adjustable model parameters &
      How should each parameter change to reduce the loss? \\
    Activations & Introduce nonlinear transformations between layers &
      Which units are active and how do they transmit gradients? \\
    Loss function & Converts predictions and targets into one scalar &
      What does ``better'' mean for this task? \\
    Optimizer & Turns gradients into parameter updates &
      How large and in which direction should the next step be? \\
    \bottomrule
  \end{tabular}
  \caption{Training a neural network requires both a model architecture and a
  procedure for turning prediction errors into parameter updates.}
\end{table}

\subsection{Backpropagation}

Backpropagation is an efficient application of the chain rule to a computational
graph. The forward pass computes the network output and stores intermediate
quantities. The backward pass starts from the derivative of the loss with
respect to the output and propagates sensitivity information back through the
same operations in reverse order.

\paragraph{A graph worked by hand.}
The mechanism is clearest on a tiny scalar graph, $f = (x+y)\cdot z$ with
$x=-2$, $y=5$, $z=-4$ (Figure~\ref{fig:backprop-graph}). The forward pass
computes and stores the intermediate value $q = x + y = 3$ and the output
$f = qz = -12$. The backward pass starts with $\partial f/\partial f = 1$ and
applies one local rule per node. The multiplication node sends to each input the
upstream gradient times the \emph{other} input:
$\partial f/\partial z = q = 3$ and $\partial f/\partial q = z = -4$. The
addition node has local derivative $1$ for both inputs, so it passes the
upstream gradient unchanged to each of them:
$\partial f/\partial x = \partial f/\partial y = -4$.

\begin{figure}[ht]
  \centering
  \begin{tikzpicture}[
      node distance=1.9cm,
      var/.style={draw, circle, minimum size=0.75cm, inner sep=0pt},
      op/.style={draw, rounded corners, fill=blue!8, minimum size=0.75cm},
      arrow/.style={-{Latex[length=2.4mm]}, thick}
    ]
    \node[var] (x) {$x$};
    \node[var, below=0.9cm of x] (y) {$y$};
    \node[op, right=of $(x)!0.5!(y)$] (add) {$+$};
    \node[var, below=0.9cm of y] (z) {$z$};
    \node[op, right=of add, yshift=-0.9cm] (mul) {$\times$};
    \node[var, right=1.6cm of mul] (f) {$f$};
    \draw[arrow] (x) -- (add);
    \draw[arrow] (y) -- (add);
    \draw[arrow] (add) -- (mul);
    \draw[arrow] (z) -- (mul);
    \draw[arrow] (mul) -- (f);
    \node[left=0.1cm of x, font=\small] {$-2$ \textcolor{red!70!black}{$(-4)$}};
    \node[left=0.1cm of y, font=\small] {$5$ \textcolor{red!70!black}{$(-4)$}};
    \node[left=0.1cm of z, font=\small] {$-4$ \textcolor{red!70!black}{$(3)$}};
    \node[above=0.3cm of $(add)!0.5!(mul)$, font=\small]
      {$q=3$ \textcolor{red!70!black}{$(-4)$}};
    \node[right=0.1cm of f, font=\small] {$-12$ \textcolor{red!70!black}{$(1)$}};
  \end{tikzpicture}
  \caption{Backpropagation on $f=(x+y)z$. Black numbers are forward values,
  red numbers in parentheses are the gradients $\partial f/\partial(\cdot)$ from
  the backward pass. The multiplication node swaps its inputs; the addition node copies the
  upstream gradient to both inputs.}
  \label{fig:backprop-graph}
\end{figure}

\noindent
These local rules are worth remembering as patterns, summarized in
Table~\ref{tab:gate-rules}. In particular, an addition node does \emph{not} send
its inputs backward; it distributes the upstream gradient.

\begin{table}[ht]
  \centering
  \begin{tabular}{@{}p{2.4cm}p{3.6cm}p{7.2cm}@{}}
    \toprule
    Node & Local derivative & Backward behavior \\
    \midrule
    Addition $a+b$ & $1$ for each input & distributes: each input receives the
      upstream gradient unchanged \\
    Multiplication $ab$ & $b$ for $a$, $a$ for $b$ & swaps: each input receives
      the upstream gradient times the other input's forward value \\
    Maximum $\max(a,b)$ & $1$ for the larger, $0$ else & routes: only the larger
      input receives the gradient (ReLU is $\max(0,z)$) \\
    \bottomrule
  \end{tabular}
  \caption{Backward rules of the elementary nodes. The multiplication rule is
  why forward values must be stored until the backward pass.}
  \label{tab:gate-rules}
\end{table}

\paragraph{The layer view.}
It is helpful to name the quantities at one layer before writing the gradient.
Let $h_k$ be the activation vector that enters layer $k$. The trainable
parameters of this layer are a weight matrix $\Omega_k$ and a bias vector
$\beta_k$. The affine part of the layer first forms the pre-activation vector
\begin{equation}
  f_k = \beta_k + \Omega_k h_k,
\end{equation}
\noindent
and a nonlinearity would then produce the next activation, for example
$h_{k+1}=\phi(f_k)$. With this notation, the gradient with respect to the weight
matrix has the structure
\begin{equation}
  \frac{\partial \ell_i}{\partial \Omega_k}
  =
  \frac{\partial \ell_i}{\partial f_k}\,h_k^\top .
\end{equation}
\noindent
This is the multiplication rule at matrix scale: the effect of a weight depends
both on the upstream error signal and on the activation feeding into that
weight. If the previous activation is doubled, the same small weight change has
twice the local effect.

\begin{figure}[ht]
  \centering
  \begin{tikzpicture}[
      node distance=1.15cm,
      box/.style={draw, rounded corners, align=center, minimum width=2.2cm,
        minimum height=0.8cm},
      arrow/.style={-{Latex[length=2.6mm]}, thick}
    ]
    \node[box, fill=green!8] (x) {$x$};
    \node[box, fill=blue!8, right=of x] (h1) {$h_1$};
    \node[box, fill=blue!8, right=of h1] (h2) {$h_2$};
    \node[box, fill=orange!12, right=of h2] (yhat) {$f_\theta(x)$};
    \node[box, fill=red!8, right=of yhat] (loss) {$\ell$};

    \draw[arrow] (x) -- (h1);
    \draw[arrow] (h1) -- (h2);
    \draw[arrow] (h2) -- (yhat);
    \draw[arrow] (yhat) -- (loss);

    \draw[arrow, dashed] (loss.south) to[out=-90,in=-90]
      node[midway, below=0.35cm, fill=white, inner sep=1pt]{gradients}
      (yhat.south);
    \draw[arrow, dashed] (yhat.south) to[out=-90,in=-90] (h2.south);
    \draw[arrow, dashed] (h2.south) to[out=-90,in=-90] (h1.south);
  \end{tikzpicture}
  \caption{The forward pass evaluates the network from left to right. The
  backward pass propagates derivatives from the loss back through the stored
  intermediate quantities.}
\end{figure}

\noindent
For ReLU hidden units, the backward recurrence includes the pointwise derivative
of the activation:
\begin{equation}
  \frac{\partial \ell_i}{\partial f_{k-1}}
  =
  \mathbb{I}[f_{k-1} > 0]
  \odot
  \left(
    \Omega_k^\top
    \frac{\partial \ell_i}{\partial f_k}
  \right),
\end{equation}
\noindent
where $\odot$ denotes elementwise multiplication. The indicator gates the
gradient: inactive ReLU units do not pass local gradient information backward.

\paragraph{What backpropagation costs and what it does not do.}
The important point is not to memorize every indexed expression, but to know
what must be stored and what is reused. One parameter update needs exactly
\emph{one} forward pass (to compute and store all intermediate values) and
\emph{one} backward pass (to reuse them); every derivative is computed once and
shared by all parameters upstream of it. Because the forward values of every
node are kept until the backward pass, the memory consumption grows with the
combined output dimensionality of all nodes in the graph (the stored
activations, times the batch size), not merely with the number of trainable
parameters. Backpropagation is only a way to \emph{compute} the gradient
exactly and efficiently. It does not change which gradient is computed, so it
neither reduces the noise of stochastic gradient descent (that comes from using
mini-batches) nor guarantees that the optimizer reaches the optimal solution.

\subsection{Automatic Differentiation and Practical Training}

Modern deep-learning libraries implement backpropagation through automatic
differentiation. A model is represented as an acyclic graph of operations. Each
operation knows how to compute its output in the forward direction and how to
propagate derivatives in the reverse direction, exactly as the nodes of
Table~\ref{tab:gate-rules}. The user specifies the model, the loss, and the
optimizer; the framework records the graph and applies the chain rule. Branches
are allowed as long as the graph stays acyclic.

A typical training loop has the following structure:
\begin{enumerate}
  \item initialize the network parameters;
  \item choose a loss function and optimizer;
  \item load a batch of examples;
  \item run the forward pass and compute the loss;
  \item run the backward pass to compute gradients;
  \item update the parameters and repeat.
\end{enumerate}
\noindent
This is the conceptual core behind PyTorch-style code: zero old gradients,
evaluate the model, call backward on the loss, and let the optimizer step.
Poor initialization can make training slow or unstable. Too large a learning
rate can overshoot; too small a learning rate can make progress painfully slow.

\subsection{Generalization: Overfitting, Bias, Variance, and Regularization}

A network that fits its training data perfectly may still be useless. The goal
is low error on \emph{new} data, and the training error cannot measure that:
a flexible model can memorize the training set, noise included, which is called
\emph{overfitting}. The opposite failure, \emph{underfitting}, occurs when the
model is not complex enough to match the data and performs poorly even on the
training set. Figure~\ref{fig:overfit} shows the typical picture: as model
complexity grows, the training error keeps falling, while the error on held-out
data first falls and then rises again once the model starts fitting noise.

\begin{figure}[ht]
  \centering
  \begin{tikzpicture}[x=1cm, y=1cm]
    \draw[-{Latex[length=2.4mm]}, thick] (0,0) -- (7.2,0)
      node[right, font=\small, align=left] {model\\complexity};
    \draw[-{Latex[length=2.4mm]}, thick] (0,0) -- (0,3.6)
      node[above, font=\small] {error};
    \draw[very thick, blue!70!black] (0.3,3.2) .. controls (2,0.9) and (4,0.4)
      .. (7,0.25);
    \draw[very thick, red!70!black] (0.3,3.3) .. controls (2,0.9) and (3.2,0.7)
      .. (4,1.0) .. controls (5,1.5) and (6,2.3) .. (7,2.9);
    \node[font=\small, text=blue!70!black, anchor=west] at (5.3,0.6)
      {training};
    \node[font=\small, text=red!70!black, anchor=west] at (5.9,3.1)
      {validation};
    \draw[dashed, gray] (3.05,0) -- (3.05,3.0);
    \node[font=\small, align=center] at (1.4,-0.45) {underfitting\\high bias};
    \node[font=\small, align=center] at (5.3,-0.45) {overfitting\\high variance};
  \end{tikzpicture}
  \vspace{0.3cm}
  \caption{Training and validation error versus model complexity. The dashed
  line marks the complexity with the lowest validation error; to its left the
  model underfits, to its right it overfits.}
  \label{fig:overfit}
\end{figure}

\noindent
The \emph{bias--variance} vocabulary names the two failures. A model with high
\emph{bias} does not match the training data closely; it underfits.
\emph{Variance} describes how much the fitted model changes when it is trained on
different portions of the data. It can be measured exactly that way: fit the
model to several subsets and compare the results. A high-variance model follows
the particular sample it saw and therefore overfits.

\paragraph{Training, validation, and test data.}
Because the training error is optimistically biased, the available data are
split into three parts with different roles. The \emph{training set} fits the
parameters. The \emph{validation set} is used for every decision about the
model: hyperparameters such as learning rate, weight decay, or depth, and when
to stop training. The degree of overfitting is read from the gap between
training and validation error, never from the training data alone. The
\emph{test set} is touched once, at the end, for an unbiased estimate; once it
influences a decision, it is no longer a test set. This separation is needed for
every model, convex or not, because convexity concerns only how well the
optimizer finds the minimum of the training loss, not how well that minimum
generalizes. It also fixes the reading of small errors: a small training error
only suggests that the model has sufficient complexity, while a small
\emph{validation} error suggests good performance.

\paragraph{Regularization.}
Regularization is any change that reduces overfitting by encoding a preference
for simpler or more robust solutions. The standard tools are:
\begin{itemize}
  \item \emph{Weight decay} ($L_2$ penalty): add $\lambda\lVert W\rVert_2^2$ to the
        loss so that large weights are discouraged.
  \item \emph{Dropout}: randomly set a fraction of activations to zero during
        training so that the network cannot rely on any single path; at test
        time all units are used.
  \item \emph{Data augmentation}: enlarge the training set with transformed
        copies (crops, flips, small rotations, color jitter). This has a
        regularizing effect, because the model must produce the same answer for
        all variants; the transformations must respect the semantics of the
        task.
  \item \emph{Early stopping}: monitor the validation loss after every epoch,
        stop when it has not improved for a fixed number $p$ of epochs
        (patience), and keep the checkpoint with the best validation loss, not
        the last one.
  \item \emph{Batch normalization}: normalize activations over the mini-batch and
        learn a scale and shift; mainly a stabilizer of optimization with a mild
        regularizing side effect.
\end{itemize}
\noindent
More data and a smaller model act in the same direction. Table~\ref{tab:diagnosis}
turns these ideas into a diagnosis checklist.

\begin{table}[ht]
  \centering
  \begin{tabular}{@{}p{3.0cm}p{5.2cm}p{5.0cm}@{}}
    \toprule
    Symptom & Meaning & Typical response \\
    \midrule
    High training error & The model or optimization setup is too weak (high
      bias) & increase capacity, train longer, adjust learning rate \\
    Low training error but high validation error & The model overfits the
      training set (high variance) & use regularization, more data, or early
      stopping \\
    Unstable loss & Updates are too aggressive or scales are poor &
      lower learning rate, improve initialization, normalize inputs \\
    \bottomrule
  \end{tabular}
  \caption{Most neural-network training failures can be diagnosed by comparing
  training behavior with validation behavior.}
  \label{tab:diagnosis}
\end{table}
